%% LyX 2.4.4 created this file.  For more info, see https://www.lyx.org/.
%% Do not edit unless you really know what you are doing.
\documentclass[english]{article}
\usepackage[T1]{fontenc}
\usepackage[utf8]{inputenc}
\setcounter{secnumdepth}{-2}
\setlength{\parindent}{0bp}
\usepackage{amsmath}
\usepackage{amsthm}
\usepackage{amssymb}

\makeatletter
%%%%%%%%%%%%%%%%%%%%%%%%%%%%%% Textclass specific LaTeX commands.
\theoremstyle{remark}
\newtheorem*{rem*}{\protect\remarkname}
\theoremstyle{plain}
\newtheorem*{lem*}{\protect\lemmaname}
\theoremstyle{plain}
\newtheorem*{thm*}{\protect\theoremname}

\@ifundefined{date}{}{\date{}}
%%%%%%%%%%%%%%%%%%%%%%%%%%%%%% User specified LaTeX commands.
\usepackage{graphicx}
\usepackage{geometry}
\newcommand{\limi}{\underset{n\rightarrow\infty}{\lim}}
\newcommand{\limxz}{\underset{x\rightarrow x_{0}}{\lim}}
\newcommand{\limfxz}{\underset{x\rightarrow x_{0}}{\lim}f(x)}
\newcommand{\limfxm}{\underset{x\rightarrow x_{0}^{-}}{\lim}f(x)}
\newcommand{\limfxp}{\underset{x\rightarrow x_{0}^{+}}{\lim}f(x)}
\newcommand{\limxm}{\underset{x\rightarrow x_{0}^{-}}{\lim}}
\newcommand{\limxp}{\underset{x\rightarrow x_{0}^{+}}{\lim}}
\newcommand{\sqrm}{M_{m\times m}(\mathbb{F})}
\newcommand{\seqa}{(a_{n})_{n=1}^{\infty}}
\newcommand{\seqx}{(x_{n})_{n=1}^{\infty}}
\newcommand{\funcdr}{f:D\rightarrow\mathbb{R}}
\newcommand{\der}{\limxz\frac{f(x)-f(x_{0})}{x-x_{0}}}
\usepackage[all]{pdfprivacy}

\makeatother

\usepackage{babel}
\providecommand{\lemmaname}{Lemma}
\providecommand{\remarkname}{Remark}
\providecommand{\theoremname}{Theorem}

\begin{document}
\title{{\Large The Bolzano-Weierstrass Theorem\vspace{-2em}}}
\maketitle
\begin{rem*}
When I write monotonically increasing or decreasing, I refer to the
weak sense.

For the strong sense, I add the word ``strictly''. 
\end{rem*}
\begin{lem*}
Every sequence $\left(a_{n}\right)_{n=1}^{\infty}$ in $\mathbb{R}$
has a monotone subsequence (a subsequence that is either non-decreasing
or non-increasing).
\end{lem*}
\begin{proof}
Let $\left(a_{n}\right)_{n=1}^{\infty}$ be a sequence in $\mathbb{R}$. 

Let us look at 
\[
A=\left\{ m\in\mathbb{N}:\forall n>m,a_{m}\ge a_{n}\right\} 
\]

This is commonly called the set of peaks of $\left(a_{n}\right)$.

\textbf{Case 1:} $A=\emptyset$ or $A$ has an upper bound.

If $A=\emptyset$, we will take $j=1$. Otherwise, we will assume
$A\ne\emptyset$ and that $A$ is bounded from above.

Let $M\in\mathbb{R}$ be an upper bound of $A$, thus, the Archimedean
property allows us to take a $j\in\mathbb{N}$ s.t. $j>M$.

Thus, every index that is bigger or equal to $j$ is not in $A$,
hence

\[
\forall n\ge j,\ \exists n'>n,\ s.t.\ a_{n}<a_{n'}
\]

We shall inductively construct a subsequence. Our first index will
be $n_{1}=j$, and for every index $j\le n_{k}$ the formula above
provides a $n_{k}'>n_{k}\ge j$ such that $a_{n_{k}}<a_{n_{k}'}$,
so we shall define $n_{k+1}$ as this $n_{k}'$.

Then $\left(n_{k}\right)_{k=1}^{\infty}$ is a strictly monotonically
increasing sequence of natural numbers, which means it successfully
defines a subsequence of $\left(a_{n}\right)$, and again from our
formula $\forall k\in\mathbb{N},\ a_{n_{k}}<a_{n_{k+1}}$ which means
this subsequence is strictly monotonically increasing, thus monotonically
increasing in particular, as required.\\

\textbf{Case 2:} Otherwise, $A$ has no upper bound and isn't empty.

This means that for every index $m\in A$, there exists an index $w>m$
s.t. $w\in A$.

We shall inductively construct a subsequence. 

As a set defined by indexes and by our assumption, $\emptyset\ne A\subseteq\mathbb{N}$,
so by the well-ordering principle of $\mathbb{N}$ it has a minimum.

Our first index will be $n_{1}=\min\left(A\right)$, and for every
index $n_{k}$ such that $n_{k}\in A$, we know there is another index
$p>n_{k}$ s.t. $p\in A$.

We shall define $n_{k+1}=p$.

Then $\left(n_{k}\right)_{k=1}^{\infty}$ is a strictly monotonically
increasing sequence of natural numbers, which means it successfully
defines a subsequence of $\left(a_{n}\right)$, and by the definition
of $A$,

\[
\forall k\in\mathbb{N},\quad a_{n_{k+1}}\le a_{n_{k}}
\]

thus $\left(a_{n_{k}}\right)_{k=1}^{\infty}$ is monotonically decreasing,
as needed. 
\end{proof}
\begin{thm*}[The Bolzano-Weierstrass theorem]
Every bounded sequence $\left(a_{n}\right)_{n=1}^{\infty}$ in $\mathbb{R}$
has a convergent subsequence.
\end{thm*}
\begin{proof}
Let $\left(a_{n}\right)_{n=1}^{\infty}$ be a bounded sequence in
$\mathbb{R}$. 

By the lemma we proved, there exists a subsequence $\left(a_{n_{k}}\right)_{k=1}^{\infty}$
that is monotone, and it is also bounded since boundedness is inherited
by subsequences. Let us recall that a sequence that is both monotone
and is bounded converges (the Monotone Convergence Theorem), therefore
$\left(a_{n_{k}}\right)_{k=1}^{\infty}$ converges and our proof is
finished.
\end{proof}

\end{document}
